\section{Posterior fields along the tree}
\label{sec:observations}

The observation construction reduces the factor problem to a question
about conditional spin laws.  We first describe those laws on a finite
subtree, where the Gaussian likelihood is explicit.  The tree structure
then expresses each covariance as a product along the connecting path.
Passing these formulas to the whole tree will identify the response
coefficients whose squared rows must be integrable in time.

For the affirmative implication, write \(b=d-1\) and assume
\(0<\theta=\tanh\beta\le b^{-1/2}\); the case \(\theta=0\) is already
IID.  On an auxiliary probability space sample the free Ising
configuration \(\sigma\), and independently sample standard Brownian
motions \((B_v)_{v\in\mathcal V}\), independent across vertices.  Set
\begin{equation}
\label{eq:observations}
 X_v(t)=t\sigma_v+B_v(t),\qquad
 \mathcal F_t^0
   =\sigma\{X_v(s):v\in\mathcal V,\ 0\le s\le t\}.
\end{equation}
The \(B_v\)'s are the noise used to plant these observations.  The
innovations \(W_v\), obtained later by subtracting the conditional drift
from \(X_v\), are different processes.  In this section
\(\mathcal F_t^0\) is the raw observation filtration; we reserve
\(\mathcal F_t\) for its usual augmentation.  For any
\(D\subseteq\mathcal V\), write
\[
 \mathcal F_t^D
   =\sigma\{X_u(s):u\in D,\ 0\le s\le t\}.
\]

The Gaussian observation framework follows Nam--Sly--Zhang
\cite{NSZ}.  In finite volume it is the observation form of stochastic
localization described by El Alaoui--Montanari
\cite[Section~3]{ElAlaouiMontanari2022}.  We derive the finite posterior
directly, since its derivative and its tree factorization will both be
needed.

\subsection{The finite posterior and its branch messages}

Let \(D\) be a finite connected subtree.  Its broadcast marginal is
\[
 \P(\sigma_D=\eta)
   =2^{-|D|}\prod_{\{i,j\}\in\mathcal E(D)}
                (1+\theta\eta_i\eta_j)
   \ \propto\
     \exp\left(\beta\sum_{\{i,j\}\in\mathcal E(D)}
                              \eta_i\eta_j\right),
 \qquad \eta\in\{-1,1\}^D.
\]
Indeed, choose any vertex of \(D\) as the broadcast root and multiply
its probability \(1/2\) by the transition probabilities on the
\(|D|-1\) edges; also
\(1+\theta s=e^{\beta s}/\cosh\beta\) for \(s\in\{-1,1\}\).
Conditional on \(\sigma_D=\eta\), the likelihood of the observation
paths through time \(t\), relative to product Wiener measure, is
\[
 \exp\left(\sum_{u\in D}\eta_uX_u(t)-\frac{t|D|}{2}\right).
\]
The last term is independent of \(\eta\).  Bayes' formula gives, for
every finite terminal vector \(x\in\mathbb R^D\),
\begin{equation}
\label{eq:finite-posterior}
 \pi_x^D(\eta)
  =\frac{1}{Z_D(x)}
     \exp\left(
       \beta\sum_{\{i,j\}\in\mathcal E(D)}\eta_i\eta_j
       +\sum_{u\in D}x_u\eta_u\right),
 \qquad x=X_D(t).
\end{equation}
Here \(Z_D(x)\) is the positive finite partition function.  This
calculation shows that the terminal vector is sufficient for the spin
posterior; it does not identify the terminal-vector sigma-field with
the entire observation history.

Let
\[
 M_v^D(x)=\sum_\eta\eta_v\pi_x^D(\eta),\qquad
 m_v^D(t)=M_v^D(X_D(t)),\qquad
 h_v^D(t)=\operatorname{atanh}m_v^D(t).
\]
By Bayes' formula, \(m_v^D(t)\) is a version of
\(\E[\sigma_v\mid\mathcal F_t^D]\).  All posterior masses are positive,
so \(|M_v^D(x)|<1\) and the field is finite.  The functions \(M_v^D\)
are smooth.  If \(u,v\in D\), then
differentiating the finite sum in \eqref{eq:finite-posterior} gives
\begin{equation}
\label{eq:finite-posterior-derivative}
 \frac{\partial M_v^D}{\partial x_u}(x)
   =\Cov_{\pi_x^D}(\eta_v,\eta_u).
\end{equation}
At \(x=X_D(t)\) we denote this conditional covariance by \(C_{vu}^D(t)\).

For an oriented edge \(v\to w\), let \(T_{v\to w}\) be the component
containing \(v\) after deletion of \(\{v,w\}\).  If both vertices belong
to \(D\), define
\[
 h_{v\to w}^D(t)=h_v^{D\cap T_{v\to w}}(t).
\]
Thus a component field uses only its own component's observations and
the unconditional free prior on that component.  It is not the field at
\(v\) after observing the opposite side or specifying the spin at \(w\).
The same distinction will be essential on the infinite tree.

\begin{lemma}[Posterior identities]
\label{lem:posterior}
For a fixed deterministic time \(t\ge0\), set
\[
 m_v(t)=\E[\sigma_v\mid\mathcal F_t^0],\qquad
 m_{v\to w}(t)
   =\E[\sigma_v\mid\mathcal F_t^{T_{v\to w}}].
\]
These means lie in \((-1,1)\) almost surely.  Write
\(h_v=\operatorname{atanh}m_v\) and
\(h_{v\to w}=\operatorname{atanh}m_{v\to w}\), and define
\[
 \psi(z)=\operatorname{atanh}(\theta\tanh z),\qquad
 a(z)=\frac{1-\tanh^2z}{1-\theta^2\tanh^2z}.
\]
At each fixed time, almost surely,
\begin{equation}
\label{eq:cavity}
 \begin{aligned}
 h_{v\to w}(t)
   &=X_v(t)+\sum_{\substack{u\sim v\\u\ne w}}
                       \psi(h_{u\to v}(t)),\\
 h_v(t)
   &=X_v(t)+\sum_{u\sim v}\psi(h_{u\to v}(t))
     =h_{v\to w}(t)+\psi(h_{w\to v}(t)).
 \end{aligned}
\end{equation}
The fields have causal, progressively measurable Borel versions as
functions of the observation paths.  The component version uses only
paths in \(T_{v\to w}\).  These versions require no distinguished
root and satisfy, at every input and time,
\begin{equation}
\label{eq:site-field-bound}
 |h_v(t)-X_v(t)|\le d\beta,\qquad
 |h_{v\to w}(t)-X_v(t)|\le b\beta.
\end{equation}
In the process formulas below, use these chosen field versions and set
\(m_v(t)=\tanh h_v(t)\) and
\(m_{v\to w}(t)=\tanh h_{v\to w}(t)\).
For the simple path \(v=v_0,v_1,\ldots,v_r=u\), define
\begin{equation}
\label{eq:covariance}
 \begin{aligned}
 c_{vu}(t)
   &=\theta^r\prod_{i=1}^r
                  a(h_{v_i\to v_{i-1}}(t)),&
 c_{vv}(t)&=1,\\
 C_{vu}(t)&=(1-m_v(t)^2)c_{vu}(t).
 \end{aligned}
\end{equation}
The last expression is a progressively measurable version of
\(\Cov(\sigma_v,\sigma_u\mid\mathcal F_t^0)\) at each fixed time.
The cavity and covariance formulas also hold in every finite connected
subtree, with branches restricted to that subtree.  For all real \(z\),
\[
 0<a(z)\le1,\qquad |\psi(z)|\le\beta,\qquad
 \psi'(z)=\theta a(z)\in(0,\theta],
\]
and \(0\le C_{vu}(t)\le\theta^{\dist(v,u)}\) at each fixed time.
\end{lemma}

\begin{proof}
We first establish the identities in a finite connected subtree \(D\).
A branch whose root has field \(z\) assigns that root probabilities
proportional to \(e^{z}\) and \(e^{-z}\).  Summing its root spin against
the edge interaction gives a factor proportional to
\(\cosh(z+\beta s)\) when its parent spin is \(s\).  Its contribution
to the parent's half log-odds is therefore
\[
 \frac12\log\frac{\cosh(z+\beta)}{\cosh(z-\beta)}
   =\operatorname{atanh}(\theta\tanh z)=\psi(z).
\]
The branch likelihoods multiply because the components are disjoint
conditional on the parent spin.  Adding their half log-odds and the
own observation \(X_v(t)\) proves the finite cavity formulas.
Since \(|\psi|\le\beta\), the site field differs from \(X_v(t)\) by at
most \(d\beta\), and a component field by at most \(b\beta\).

We next follow a path \(v_0,\ldots,v_r\) in \(D\).  With the observations
fixed, sum the posterior in \eqref{eq:finite-posterior} over all spins
off the path.  Given the preceding path spins, the conditional mean
of the spin at \(v_i\) uses only \(\sigma_{v_{i-1}}\) and the field on
the side excluding that predecessor:
\[
 \E[\sigma_{v_i}\mid \sigma_{v_0},\ldots,\sigma_{v_{i-1}},
                    \mathcal F_t^D]
   =\tanh(z+\beta\sigma_{v_{i-1}})
   =p_i+q_i\sigma_{v_{i-1}},
 \qquad z=h_{v_i\to v_{i-1}}^D(t).
\]
Here \(p_i\) is an intercept depending on the observations, and
\[
 q_i=\frac{\tanh(z+\beta)-\tanh(z-\beta)}{2}
     =\theta\,\frac{1-\tanh^2z}{1-\theta^2\tanh^2z}
     =\theta a(z).
\]
Iterating these affine conditional means makes the slope of the
conditional mean at \(v_r\), as a function of \(\sigma_{v_0}\),
equal to \(\prod_{i=1}^r q_i\).  Taking covariance with the first
spin gives
\[
 C_{v_0v_r}^D(t)
   =(1-(m_{v_0}^D(t))^2)\prod_{i=1}^r q_i.
\]
This proves the finite path formula, including the empty product at
\(r=0\).  The coefficient \(c_{vu}^D\) is directed even though the
covariance is symmetric: its normalization uses the variance at \(v\).
In fact, \eqref{eq:finite-posterior-derivative} gives
\[
 \frac{\partial}{\partial x_u}\operatorname{atanh}M_v^D(x)
   =\frac{\Cov_{\pi_x^D}(\eta_v,\eta_u)}
           {1-M_v^D(x)^2},
\]
which is the path product just obtained.  Thus \(c_{vu}^D\) is the
finite field's response to the observation at \(u\).

Now fix \(t\) and exhaust the tree by increasing finite balls.
The conditional expectations of a fixed spin, or a product of two
fixed spins, converge almost surely and in \(L^2\) to conditioning
on \(\mathcal F_t^0\).  This is increasing-sigma-field convergence
for bounded random variables.  Intersecting the balls with
\(T_{v\to w}\) gives the corresponding convergence for the component
posterior.  The finite field bounds place each mean in the random
compact interval
\[
 [\tanh(X_v(t)-k\beta),\,\tanh(X_v(t)+k\beta)]
       \ \subset\ (-1,1),\qquad k=d\ \hbox{or}\ b.
\]
Thus the limiting means remain strictly inside \((-1,1)\), and their
inverse hyperbolic tangents converge to finite fields.
For a fixed path, take balls large enough to contain it.  Their
intersections with the finitely many directed components on the path
exhaust those components.  The finite cavity sums and path product
therefore pass to the limit; convergence of the two-spin moments
identifies the limit as the conditional covariance.  Any other centered
exhaustion gives the same conditional expectations, since its histories
generate the same sigma-fields.

It remains to choose process versions without presupposing a continuous
infinite-dimensional filter.  On the canonical space of continuous
observation paths, evaluate the finite field formulas on \(B_R(v)\)
for a site and on \(B_R(v)\cap T_{v\to w}\) for a component, with
\(R=0,1,\ldots\).  Define each raw field to be the pointwise
\(\limsup_{R\to\infty}\) of these finite fields.  Each finite expression
is a continuous adapted terminal-array function.  Its upper limit is
Borel and progressively measurable, for the whole observation
filtration or for the specified component filtration respectively.
The finite bounds make the upper limits finite and preserve
\eqref{eq:site-field-bound} at every input and time.  At each fixed
time they agree almost surely with the conditional fields proved above.
Their hyperbolic tangents and the finite products in
\eqref{eq:covariance} give the claimed mean and covariance versions.

These formulas use only the terminal array at the evaluated time, and
in particular only the histories up to that time.  Centering every
ball at its indexed vertex makes the formulas commute with graph
relabeling.  Joint measurability and Fubini's theorem turn the
fixed-time algebraic identities into identities for almost every time,
almost surely.  Countability of the vertices, oriented edges, finite
paths, and integer time horizons permits all these time-probability
statements to hold simultaneously.  Continuity and identities
simultaneous at every time are not needed here; they will be obtained
after the filtering argument in Section~\ref{sec:innovations}.

Finally, direct differentiation gives
\(\psi'(z)=\theta(1-\tanh^2z)/(1-\theta^2\tanh^2z)\).
Since \(0<\theta<1\), its bounds and \(0<a\le1\) follow.
The covariance bound follows from the finite product and
\(1-m_v^2\le1\).
\end{proof}

The normalized response \(c_{vu}\), rather than the covariance alone,
is the coefficient that will occur in the field equation.  At distance
\(r\), its deterministic bound is \(\theta^r\).  Below the threshold
this already gives a summable squared row.  At equality, the size of a
distance shell cancels \(\theta^{2r}\); the next section obtains the
additional decay from the observations in side branches.
